\section{Data appendix}\label{app:data}

\subsection{Sample and variable construction}

I pool the twenty quarterly ENAHO employment modules (module 500) for 2021Q1--2025Q4;
the main analysis uses 2021Q1--2024Q4 and the validation uses 2024Q1--2025Q4. The
working-age sample comprises individuals aged 14 to 65 with non-missing education. All
statistics use the module employment weight (\texttt{fac500a}). This weight is
calibrated so that the full year of interviews expands to the annual population
projection; each quarterly subsample therefore expands to roughly one fourth of the
population, which is immaterial for the weighted means, shares, and regression
estimates used throughout (numerator and denominator scale together) and would matter
only for population totals, which I never estimate. The quarterly partition follows
the interview month and matches the survey's own design: quarterly national inference
is an official inference level of ENAHO in the technical sheets of both 2021 and 2025,
and the non-panel sample revisits the same conglomerates in the same months every
year. I verify both properties in the microdata for the full window: all 5{,}359
conglomerates appear in every year from 2021 to 2025, the overlap of primary sampling
units between the same quarter of consecutive years ranges from 70 to 92 percent
(against 2 to 5 percent between different quarters), and each quarter is an
internally balanced replica of the design. Two survey changes over the window are
worth recording. First, the mixed interview modality adopted during the pandemic
(face-to-face prioritized, telephone follow-up for absent members) remained the
standard through 2025, so the mode is broadly consistent across my study period; the
one exception is February 2021, when a reduced telephone questionnaire in fifteen
regions dropped, among other items, the contract-type question (missing for half of
that month's wage earners), which affects none of the baseline outcomes (2021
informality comes from INEI's official indicator and income items were retained) and
is handled explicitly in the one appendix margin that uses contracts. Second, the
2025 technical sheet describes the design as two-stage with urban conglomerates of
about 140 dwellings, against the multi-stage description of 2021; the conglomerate
identifiers and counts are nonetheless continuous across the window, so this is a
reframing of the documentation rather than a break in the sample. Departments are the 25
first-level administrative units derived from the first two digits of \texttt{ubigeo};
the urban indicator is derived from the sampling stratum.

\paragraph{Education and skill groups.} Educational attainment is ENAHO variable
\texttt{p301a}. I verified its coding against the weighted attainment distribution.
Low-skilled workers are those with \texttt{p301a} in 1--6 (no education through complete
secondary); high-skilled workers are those with \texttt{p301a} in 7--11 (non-university
higher through postgraduate). Special education (code 12) and missing values are excluded.

\paragraph{Labor income and wages.} ENAHO records the pay of wage earners per pay period:
\texttt{p524a1} is the gross payment including overtime and \texttt{p523} its frequency
(daily, weekly, fortnightly, or monthly). Fifty-four percent of low-education wage
earners are paid at non-monthly frequencies, against fifteen percent of high-education
wage earners, so I convert all payments to monthly equivalents using the same factors
implied by INEI's own annualized income variable (\texttt{i524a1}): 25 for daily, $52/12$
for weekly, 2 for fortnightly, and 1 for monthly. The full weighted frequency
distribution among private wage earners with valid earnings is 2 percent daily, 45
weekly, 9 fortnightly, and 45 monthly for the low-education group, against 1, 18, 6,
and 76 percent for the high-education group; the frequency variable is never missing
in that sample. Wage earners comprise employees and
domestic workers (\texttt{p507} in $\{3,4,6\}$); domestic workers have been entitled to
the minimum wage since Law 31047. For employers and own-account workers, labor income is
net profit (\texttt{p530a}), which is already monthly. Because reported income refers to
the month before the interview, I deflate at that reference month (January interviews
roll back to December of the previous year) using the Lima consumer price index
rebased to December 2021, and I evaluate the statutory minimum in force in the same
reference month when classifying workers as paid below the floor. The real hourly wage
divides monthly-equivalent labor income by $4.345$ times weekly hours in the main
occupation (\texttt{p513t}; INEI's imputed hours variable \texttt{i513t} substitutes
where \texttt{p513t} is missing, but no observation in the wage-analysis sample and
fewer than 0.01 percent of the hours sample require it). Wage outcomes exclude public-sector
workers, whose pay is set by separate scales, and are winsorized once at the 0.5th and
99.5th percentiles of the pooled private wage-earner sample.

\paragraph{Informality.} INEI's official informal-employment indicator (\texttt{ocupinf},
coded 1 for informal and 2 for formal in these files) is present for 2021--2023 but absent
for 2024--2025. I reconstruct a consistent indicator from primitives available in every
year. A wage employee or domestic worker is classified as formal if affiliated to a
pension system (\texttt{p558a1}--\texttt{p558a4}, whose selected options are marked with
their option codes rather than with ones) or holding a formal labor contract
(\texttt{p511a} in 1--6); an employer or own-account worker is classified as formal if the
productive unit is registered (\texttt{p510a1}); unpaid family workers are always informal.
Validated against \texttt{ocupinf} on the overlapping years 2021--2023, this rule agrees
with the official classification for 87.8 percent of employed workers and yields an
informality rate of 74.9 percent against the official 79.1 percent. I use the official
indicator wherever available and the reconstruction only for 2024--2025; within the
2024--2025 validation window the indicator is reconstructed in every quarter, so it is
measured consistently there.

\paragraph{Minimum wage and bite.} The statutory minimum wage schedule is taken from the
decrees and verified against the Banco Central de Reserva del Perú series (PN02124PM);
the CPI deflator is BCRP series PN38705PM (Lima, December 2021 $=$ 100), and both inputs
are regenerated by a documented script in the replication package. The department Kaitz
index is the ratio of the post-reform minimum (1,025 soles) to the department's
pre-reform (2021Q1--2022Q1) weighted median monthly-equivalent wage of private wage
earners. The fraction-affected measure is the pre-reform weighted share of private wage
earners paid below 1,025 soles.

\paragraph{The ENAHO Panel.} The transition analysis of Section~\ref{sec:transitions}
uses the ENAHO Panel 2020--2024 (INEI survey 978, published November 2025), whose
employment module (\texttt{enaho01a-2020-2024-500-panel.dta}) contains one record per
person with year-suffixed variables for 2020 through 2024. Persons observed in both
years of a pair are flagged by the indicators \texttt{perpanel} (23,983 linked persons
for 2021--22 and 23,245 for 2022--23) and weighted with the INEI panel expansion
factors \texttt{facpanel}, which adjust for attrition and non-response; official
inference for consecutive-year pairs is national, urban and rural, and by natural
region. Employment states replicate the cross-section definitions exactly: employed
is \texttt{ocu500}${}=1$; the formal private state requires private-sector employment
with a formal job; informality uses the official \texttt{ocupinf} where released and
the Rule-F reconstruction above elsewhere; self-employment comprises own-account
workers and employers. Base-year education defines the skill groups. Because the
panel is annual and the reform took effect on 1 May 2022, pair-observations are
classified with the 2022 interview month as described in the notes to
Table~\ref{tab:transitions}.

\paragraph{Protest intensity.} The department-level intensity of the 2022--23
political crisis is measured from the monthly conflict reports of the Presidencia del
Consejo de Ministros (Secretar\'ia de Gesti\'on Social y Di\'alogo, ``Reporte de
Conflictos Sociales''), whose crisis-monitoring tables record the maximum number of
persons mobilized in protest actions in each department in each week and month. I take
each department's maximum across the January and February 2023 reports, the two months
in which national mobilization peaked (the December-2022 report predates the table),
divide it by the department's working-age population (four times the mean quarterly sum
of the ENAHO expansion weights), and standardize the result across the 25 departments.
The transcribed table and the source documents are included in the replication
package.

\subsection{Supplementary results}

Table~\ref{tab:ddd} reports the triple-difference estimates and Table~\ref{tab:loo}
the leave-one-out stability ranges discussed in Section~\ref{sec:robustness};
Tables~\ref{tab:expocells}--\ref{tab:margins} collect the exposure-design diagnostics,
the sectoral and occupational checks, the selection bounds and power calculations, the
synthetic-DiD and pre-test-power results, and the adjustment margins.
Table~\ref{tab:crisis} reports the political-crisis robustness of the continuous
regional-exposure design and Table~\ref{tab:uncond} the unconditional compositional
outcomes, both discussed in Section~\ref{sec:robustness}.
Table~\ref{tab:placebo} reports the in-time placebo reforms and confirms that the wage
placebos are null while the employment placebos are contaminated by the pandemic
pre-trend. The doubly robust estimates of \citet{santanna_zhao_2020} appear alongside
the two-way fixed-effects estimates in Table~\ref{tab:main}.
Figure~\ref{fig:val2025} shows the event study around the January-2025 reform used in
the out-of-sample validation.

\begin{table}[!tbp]\centering
\caption{Triple-difference estimates: $\mathrm{Low}\times\mathrm{Post}\times\mathrm{HighBite}$}\label{tab:ddd}
\begin{threeparttable}
\begin{tabular}{lc}
\toprule
Outcome & DDD estimate \\
\midrule
Log hourly wage & -0.027 \\
 & (0.038) \\
Employment & 0.015 \\
 & (0.016) \\
Formal empl. & -0.018 \\
 & (0.018) \\
Self-employment & 0.026 \\
 & (0.018) \\
\bottomrule
\end{tabular}
\begin{tablenotes}\footnotesize
\item Notes: Coefficient on the triple interaction $\mathrm{Low}\times\mathrm{Post}\times\mathrm{HighBite}$, where HighBite indicates departments with an above-median pre-reform fraction of private wage earners paid below the new minimum wage. Department and quarter fixed effects, controls as in the baseline; standard errors clustered by department (25 clusters); ENAHO survey weights. $^{*}p<0.1$, $^{**}p<0.05$, $^{***}p<0.01$.
\end{tablenotes}
\end{threeparttable}\end{table}


\begin{table}[H]\centering
\caption{Leave-one-out stability of the baseline DiD estimates}\label{tab:loo}
\begin{threeparttable}
\begin{tabular}{lccc}
\toprule
Outcome & Full sample & Drop-one-department range & Drop-one-quarter range \\
\midrule
Log hourly wage & -0.021 & [-0.027, -0.017] & [-0.026, -0.010] \\
Employment & -0.025 & [-0.028, -0.018] & [-0.031, -0.014] \\
Formal empl. & -0.046 & [-0.050, -0.043] & [-0.052, -0.042] \\
Self-employment & 0.036 & [0.032, 0.038] & [0.031, 0.041] \\
\bottomrule
\end{tabular}
\begin{tablenotes}\footnotesize
\item Notes: Ranges of the $\mathrm{Low}\times\mathrm{Post}$ coefficient when re-estimating the baseline dropping one department at a time (25 estimates) or one quarter at a time (15 estimates). Full per-unit estimates are in the replication output (\texttt{rob\_leaveoneout\_detail.csv}).
\end{tablenotes}
\end{threeparttable}\end{table}


\begin{table}[H]\centering
\caption{Exposure designs: diagnostics and cell-level estimates}\label{tab:expocells}
\begin{threeparttable}
\footnotesize
\begin{tabular}{lcccc}
\toprule
\multicolumn{5}{l}{\emph{Panel A: regional Kaitz design, exposure-leads joint tests}}\\
\midrule
Log hourly wage & \multicolumn{4}{c}{leads $p$ = 0.243} \\
Employment & \multicolumn{4}{c}{leads $p$ = 0.002} \\
Formal empl. & \multicolumn{4}{c}{leads $p$ = 0.842} \\
Self-employment & \multicolumn{4}{c}{leads $p$ = 0.202} \\
\midrule
\multicolumn{5}{l}{\emph{Panel B: cell-level fraction-affected designs (transparency exercise)}}\\
Outcome & Exposure & Estimate & (s.e.) & Leads $p$ \\ \midrule
Log hourly wage & frac\_below & 0.129$^{**}$ & (0.056) & 0.011 \\
Log hourly wage & frac\_aff & 0.016 & (0.133) & 0.004 \\
Employment & frac\_below & -0.225$^{***}$ & (0.032) & 0.000 \\
Employment & frac\_aff & 0.391$^{***}$ & (0.072) & 0.000 \\
Formal empl. & frac\_below & -0.074$^{***}$ & (0.014) & 0.007 \\
Formal empl. & frac\_aff & 0.303$^{***}$ & (0.039) & 0.031 \\
Self-employment & frac\_below & 0.079$^{***}$ & (0.018) & 0.244 \\
Self-employment & frac\_aff & -0.155$^{**}$ & (0.065) & 0.045 \\
\bottomrule
\end{tabular}
\begin{tablenotes}\footnotesize
\item Notes: Panel A: joint tests that the pre-reform quarter$\times$exposure interactions are zero in the regional continuous-exposure design (standardized department Kaitz). Panel B: coefficients on $\mathrm{Post}\times$exposure from cell-level fraction-affected designs (department$\times$skill$\times$age cells; frac\_below = pre-reform share of the cell's private wage earners paid below 1{,}025; frac\_aff = share paid between 930 and 1{,}025), with cell and quarter fixed effects and department-clustered standard errors. Both cell-level variants fail their leads diagnostics -- fine-grained exposure is confounded with the cells' pandemic-recovery trajectories -- which is why the regional design is the preferred corroboration. $^{*}p<0.1$, $^{**}p<0.05$, $^{***}p<0.01$.
\end{tablenotes}
\end{threeparttable}\end{table}


\begin{table}[!tbp]\centering
\caption{Sectoral breadth and occupational dose response of the reallocation}\label{tab:sectorocc}
\begin{threeparttable}
\footnotesize
\begin{tabular}{lccc}
\toprule
\multicolumn{4}{l}{\emph{Panel A: education DiD by sector group}}\\
Outcome & Covered high-bite & Agriculture & Other private \\
\midrule
Formal empl. & -0.029$^{***}$ & -0.051$^{***}$ & -0.044$^{***}$ \\
 & (0.006) & (0.013) & (0.012) \\
Self-employment & 0.028$^{***}$ & 0.020 & 0.030$^{**}$ \\
 & (0.007) & (0.013) & (0.012) \\
Log hourly wage & -0.040$^{**}$ & -0.060 & 0.015 \\
 & (0.018) & (0.050) & (0.014) \\
\midrule
\multicolumn{4}{l}{\emph{Panel B: occupation-exposure dose response ($\mathrm{Post}\times$ occupation share below 1{,}025)}}\\
Outcome & Estimate & (s.e.) & Leads $p$ \\ \midrule
Formal empl. & -0.049$^{***}$ & (0.016) & 0.008 \\
Self-employment & 0.000 & (0.017) & 0.036 \\
Log hourly wage & 0.131$^{***}$ & (0.022) & 0.106 \\
Paid below MW & -0.094$^{***}$ & (0.023) & 0.000 \\
\bottomrule
\end{tabular}
\begin{tablenotes}\footnotesize
\item Notes: Panel A: $\mathrm{Low}\times\mathrm{Post}$ DiD within sector groups defined from CIIU Rev.4 divisions (covered high-bite: manufacturing 10--33, construction 41--43, trade 45--47, hotels and restaurants 55--56; agriculture: divisions 01--03, under the agrarian regime of Ley 31110, whose daily remuneration is indexed to the RMV and therefore rose with the reform). Panel B: coefficient on $\mathrm{Post}\times$(pre-reform share of the 2-digit occupation's private wage earners paid below 1{,}025), with occupation, quarter, and department fixed effects; occupation is measured at the interview and is therefore post-treatment for movers, so Panel B is descriptive corroboration rather than a primary design. Department-clustered standard errors; ENAHO survey weights. $^{*}p<0.1$, $^{**}p<0.05$, $^{***}p<0.01$.
\end{tablenotes}
\end{threeparttable}\end{table}


\begin{table}[H]\centering
\caption{Selection bounds and statistical power}\label{tab:leemde}
\begin{threeparttable}
\footnotesize
\begin{tabular}{lcc}
\toprule
\multicolumn{3}{l}{\emph{Panel A: Lee-type trimming bounds, unconditional 2$\times$2 log-wage DiD}}\\
\midrule
Selection DiD (wage-sample inclusion) & -0.021 & \\
Trimming fraction & 0.076 & \\
Untrimmed DiD & -0.005 & \\
Bounds & [-0.099, 0.101] & \\
\midrule
\multicolumn{3}{l}{\emph{Panel B: minimum detectable effects (80\% power, 5\% size)}}\\
Outcome & Estimate & MDE \\ \midrule
Log real hourly wage & -0.021 & 0.034 \\
Log real monthly earnings & -0.010 & 0.038 \\
Employed (working age) & -0.025 & 0.016 \\
Labour force participation & -0.008 & 0.012 \\
Formal employment & -0.046 & 0.021 \\
Weekly hours (main job) & 0.556 & 0.698 \\
Full-time (>=35h) & -0.013 & 0.013 \\
Self-employed & 0.036 & 0.017 \\
Paid below the MW & -0.006 & 0.021 \\
\bottomrule
\end{tabular}
\begin{tablenotes}\footnotesize
\item Notes: Panel A: bounds on the unconditional 2$\times$2 log hourly wage DiD following the trimming logic of \citet{lee_2009}: the reform reduced the probability that a low-education worker appears in the wage-earner sample (selection DiD), so the low$\times$post cell is trimmed from above and below by the implied fraction of its baseline rate. Panel B: minimum detectable effect $=2.8\times$ the department-clustered standard error of the baseline DiD.
\end{tablenotes}
\end{threeparttable}\end{table}


\begin{table}[H]\centering
\caption{Synthetic DiD and the power of the pre-trends test}\label{tab:sdidpower}
\begin{threeparttable}
\footnotesize
\begin{tabular}{lccc}
\toprule
\multicolumn{4}{l}{\emph{Panel A: synthetic difference-in-differences (high-bite vs synthetic low-bite departments)}}\\
Outcome & Low-education means & Skill gap & \\
\midrule
Log hourly wage & 0.041 & -0.028 & \\
 & (0.025) & (0.048) & \\
Employment & -0.039$^{***}$ & -0.016 & \\
 & (0.010) & (0.011) & \\
Formal empl. & -0.013$^{**}$ & -0.028 & \\
 & (0.007) & (0.017) & \\
Self-employment & 0.011 & 0.023 & \\
 & (0.011) & (0.019) & \\
\midrule
\multicolumn{4}{l}{\emph{Panel B: power of the pre-trends test against linear violations \citep{roth_2022_pretest}}}\\
Outcome & Slope at 50\% power & Implied bias (avg.\ post) & Bias / estimate \\ \midrule
Log hourly wage & 0.008 & 0.052 & -2.53 \\
Employment & 0.006 & 0.036 & -1.43 \\
Formal empl. & 0.006 & 0.041 & -0.89 \\
Self-employment & 0.006 & 0.041 & 1.13 \\
\bottomrule
\end{tabular}
\begin{tablenotes}\footnotesize
\item Notes: Panel A: synthetic DiD estimates \citep{arkhangelsky_2021_sdid} on department$\times$quarter cells, treating above-median-bite departments as treated from 2022Q3; column (1) uses low-education outcome means, column (2) the low-minus-high skill gap; jackknife standard errors in parentheses (the placebo variance is infeasible with 12 treated and 13 control units). Panel B: slope of a linear violation of parallel trends against which the conventional pre-test has only 50 percent power \citep{roth_2022_pretest}, the bias such an undetected violation would induce in the average post-period effect, and its ratio to the estimated effect. $^{*}p<0.1$, $^{**}p<0.05$, $^{***}p<0.01$.
\end{tablenotes}
\end{threeparttable}\end{table}


\begin{table}[!tbp]\centering
\caption{Adjustment margins and additional subgroups}\label{tab:margins}
\begin{threeparttable}
\footnotesize
\begin{tabular}{lccc}
\toprule
\multicolumn{4}{l}{\emph{Panel A: adjustment margins ($\mathrm{Low}\times\mathrm{Post}$)}}\\
Margin & Estimate & (s.e.) & Pre-reform mean \\
\midrule
No written contract (wage earners) & 0.031$^{***}$ & (0.008) & 0.610 \\
Micro firm, <=20 workers (employed) & 0.028$^{***}$ & (0.007) & 0.854 \\
Tenure < 12 months (wage earners) & -0.012 & (0.016) & 0.513 \\
Informal job, informal sector (2021-23) & 0.024$^{**}$ & (0.010) & 0.663 \\
Informal job, FORMAL sector (2021-23) & 0.008 & (0.007) & 0.144 \\
\midrule
\multicolumn{4}{l}{\emph{Panel B: subgroups by household role and ethnic self-identification}}\\
Outcome & Group & Estimate & (s.e.) \\ \midrule
Employment & Head & -0.003 & (0.006) \\
Employment & Non-head & -0.038$^{***}$ & (0.008) \\
Employment & Indigenous & -0.004 & (0.015) \\
Employment & Non-indigenous & -0.030$^{***}$ & (0.005) \\
Formal empl. & Head & -0.005 & (0.011) \\
Formal empl. & Non-head & -0.068$^{***}$ & (0.006) \\
Formal empl. & Indigenous & -0.030 & (0.018) \\
Formal empl. & Non-indigenous & -0.048$^{***}$ & (0.008) \\
Self-employment & Head & 0.030$^{***}$ & (0.008) \\
Self-employment & Non-head & 0.039$^{***}$ & (0.008) \\
Self-employment & Indigenous & 0.009 & (0.017) \\
Self-employment & Non-indigenous & 0.045$^{***}$ & (0.008) \\
Log hourly wage & Head & -0.020 & (0.017) \\
Log hourly wage & Non-head & -0.019 & (0.014) \\
Log hourly wage & Indigenous & -0.008 & (0.029) \\
Log hourly wage & Non-indigenous & -0.026$^{*}$ & (0.013) \\
\bottomrule
\end{tabular}
\begin{tablenotes}\footnotesize
\item Notes: Panel A: $\mathrm{Low}\times\mathrm{Post}$ DiD estimates for mechanism-oriented outcomes: the share of private wage earners without a written contract (\texttt{p511a}), the share of private workers in micro units of at most 20 workers (\texttt{p512a}), the share of wage earners with under 12 months of tenure (\texttt{p513a1/2}), and the decomposition of informal employment into informal-sector and formal-sector informal jobs (INEI \texttt{emplpsec}, available 2021--2023). Panel B: baseline DiD within subgroups defined by household headship (\texttt{p203}) and indigenous self-identification (\texttt{p558c}). Department-clustered standard errors; ENAHO survey weights. $^{*}p<0.1$, $^{**}p<0.05$, $^{***}p<0.01$.
\end{tablenotes}
\end{threeparttable}\end{table}


\begin{table}[!tbp]\centering
\caption{Robustness to the Venezuelan immigration wave}\label{tab:migration}
\begin{threeparttable}
\footnotesize
\setlength{\tabcolsep}{4pt}
\begin{tabular}{lcccc}
\toprule
Specification & Log hourly wage & Employment & Formal empl. & Self-employment \\
\midrule
Baseline & -0.021 & -0.025$^{***}$ & -0.046$^{***}$ & 0.036$^{***}$ \\
 & (0.012) & (0.006) & (0.007) & (0.006) \\
Department$\times$quarter FE & -0.020 & -0.015$^{**}$ & -0.041$^{***}$ & 0.034$^{***}$ \\
 & (0.013) & (0.006) & (0.008) & (0.006) \\
Excluding Lima and Callao & -0.021 & -0.018$^{*}$ & -0.042$^{***}$ & 0.037$^{***}$ \\
 & (0.022) & (0.009) & (0.010) & (0.010) \\
Excluding top-8 migrant-hosting departments & -0.051 & -0.012 & -0.049$^{***}$ & 0.045$^{**}$ \\
 & (0.032) & (0.013) & (0.015) & (0.016) \\
\bottomrule
\end{tabular}
\begin{tablenotes}\footnotesize
\item Notes: $\mathrm{Low}\times\mathrm{Post}$ DiD estimates under specifications that address the Venezuelan immigration inflow. Department$\times$quarter fixed effects absorb any local shock common to both education groups within a department-quarter, including migrant arrivals and regularization waves. The top-8 exclusion drops Lima, Callao, La Libertad, Arequipa, Lambayeque, Piura, Ica, and Tumbes, the departments hosting the large majority of the Venezuelan population. Department-clustered standard errors; ENAHO survey weights. $^{*}p<0.1$, $^{**}p<0.05$, $^{***}p<0.01$.
\end{tablenotes}
\end{threeparttable}\end{table}


\begin{table}[!tbp]\centering
\caption{The 2022--23 political crisis and the continuous regional-exposure design}\label{tab:crisis}
\begin{threeparttable}
\footnotesize
\setlength{\tabcolsep}{3pt}
\begin{tabular}{lcccccc}
\toprule
& (1) Baseline & (2) +Post$\times$ & (3) +2023$\times$ & (4) Excl.\ south & (5) Post = & (6) (4)+(5) \\
Outcome & & protest & protest & bloc & 2022 only & \\
\midrule
Log hourly wage & 0.017 & 0.012 & 0.015 & 0.009 & 0.019 & 0.008 \\
 & (0.012) & (0.012) & (0.012) & (0.012) & (0.016) & (0.014) \\
 & [0.034] & [0.074] & [0.038] & [0.156] & [0.088] & [0.355] \\
Employment & -0.034$^{***}$ & -0.027$^{***}$ & -0.029$^{***}$ & -0.029$^{**}$ & -0.029$^{***}$ & -0.024$^{**}$ \\
 & (0.012) & (0.009) & (0.009) & (0.010) & (0.009) & (0.009) \\
 & [0.001] & [0.004] & [0.001] & [0.000] & [0.002] & [0.007] \\
Formal empl. & -0.018$^{***}$ & -0.017$^{***}$ & -0.016$^{***}$ & -0.016$^{***}$ & -0.006 & -0.003 \\
 & (0.003) & (0.003) & (0.002) & (0.003) & (0.004) & (0.003) \\
 & [0.001] & [0.001] & [0.000] & [0.002] & [0.453] & [0.877] \\
Self-employment & 0.010$^{**}$ & 0.008$^{**}$ & 0.008$^{**}$ & 0.010$^{**}$ & 0.007$^{*}$ & 0.008 \\
 & (0.004) & (0.004) & (0.003) & (0.004) & (0.004) & (0.006) \\
 & [0.023] & [0.052] & [0.044] & [0.060] & [0.069] & [0.151] \\
\bottomrule
\end{tabular}
\begin{tablenotes}\footnotesize
\item Notes: Each cell reports the coefficient on $\mathrm{Post}\times$(standardized department Kaitz index) from the continuous-exposure design, with department-clustered standard errors in parentheses and the randomization-inference $p$-value (permuting the Kaitz indices across the included departments, 1{,}000 draws, protest control held at its true assignment) in brackets. Column (1) is the baseline of Table~\ref{tab:inference}, column (5). Columns (2) and (3) add the standardized per-capita peak protest mobilization interacted with the post indicator and with an indicator for 2023Q1 onward, respectively; protest intensity is the January--February 2023 maximum number of persons mobilized in each department, from the monthly conflict reports of the Presidencia del Consejo de Ministros, divided by the department's working-age population. Column (4) excludes the six southern departments at the center of the December-2022--March-2023 protests (Apur\'imac, Arequipa, Ayacucho, Cusco, Madre de Dios, Puno). Column (5) restricts the post-reform window to 2022Q3--Q4, before the ouster of President Castillo on 7 December 2022, and column (6) combines both restrictions. All estimates weighted by ENAHO survey weights. $^{*}p<0.1$, $^{**}p<0.05$, $^{***}p<0.01$.
\end{tablenotes}
\end{threeparttable}\end{table}


\begin{table}[!tbp]\centering
\caption{Unconditional compositional outcomes (shares of the working-age population)}\label{tab:uncond}
\begin{threeparttable}
\footnotesize
\setlength{\tabcolsep}{4pt}
\begin{tabular}{lccccc}
\toprule
& Pre-reform & Education DiD & Pre-trend & Kaitz slope & RI \\
Outcome & mean & (Low$\times$Post) & $p$ & (Post$\times z$) & $p$ \\
\midrule
Private formal employment & 0.148 & -0.033$^{***}$ & 0.014 & -0.016$^{***}$ & 0.000 \\
 & & (0.006) & & (0.003) & \\
Self-employment & 0.252 & 0.017$^{***}$ & 0.000 & -0.004$^{*}$ & 0.115 \\
 & & (0.004) & & (0.003) & \\
Private informal employment & 0.500 & 0.009 & 0.306 & -0.019$^{**}$ & 0.009 \\
 & & (0.006) & & (0.008) & \\
\bottomrule
\end{tabular}
\begin{tablenotes}\footnotesize
\item Notes: Outcomes are defined for every working-age individual: an individual counts as (private) formally employed only if employed, in the private sector, and formal, and analogously for self-employment and private informal employment; the non-employed and public-sector workers count as zeros. Column (2) reports the baseline $\mathrm{Low}\times\mathrm{Post}$ education DiD with department-clustered standard errors; column (3) the joint $p$-value of the pre-reform event-study interactions; column (4) the coefficient on $\mathrm{Post}\times$(standardized department Kaitz index) from the continuous regional-exposure design; column (5) its randomization-inference $p$-value (permuting the Kaitz indices across departments, 1{,}000 draws). Because these outcomes embed the employment margin, the education-DiD versions inherit its differential pandemic-recovery pre-trend and are reported for transparency rather than leaned on; the regional-design estimates do not depend on the education contrast. All estimates weighted by ENAHO survey weights. $^{*}p<0.1$, $^{**}p<0.05$, $^{***}p<0.01$.
\end{tablenotes}
\end{threeparttable}\end{table}


\paragraph{Deflators.} All wage regressions include quarter fixed effects, which
absorb national inflation, so log-wage DiD estimates are invariant to the national
deflator; and the bite measures compare nominal wages with the nominal statutory
floor, so they involve no deflation at all. Department-specific price deflation would
matter only for cross-departmental comparisons of real wage \emph{levels}, which play
no role in the identification.

\begin{figure}[!t]
\centering
\includegraphics[width=0.9\textwidth]{fig9_validation2025.pdf}
\caption{Event study around the January-2025 reform (out-of-sample validation). Coefficients
on the interaction of the low-skill indicator with quarter indicators, reference quarter
2024Q4. Bands are ninety-five percent confidence intervals clustered by department.}
\label{fig:val2025}
\end{figure}

\begin{table}[!tbp]\centering
\caption{In-time placebo reforms (pre-reform window only)}\label{tab:placebo}
\begin{threeparttable}
\begin{tabular}{lccc}
\toprule
Outcome & Placebo 2021Q2 & Placebo 2021Q3 & Placebo 2021Q4 \\
\midrule
Log real hourly wage & 0.020 & -0.018 & -0.005 \\
 & (0.037) & (0.032) & (0.021) \\
Employment & -0.057$^{***}$ & -0.034$^{***}$ & -0.026$^{***}$ \\
 & (0.016) & (0.007) & (0.008) \\
Formal employment & -0.021$^{*}$ & -0.012 & -0.005 \\
 & (0.011) & (0.013) & (0.007) \\
Self-employment & 0.022$^{*}$ & -0.009 & 0.007 \\
 & (0.012) & (0.007) & (0.009) \\
\bottomrule
\end{tabular}

\begin{tablenotes}\footnotesize
\item Notes: Each entry is the $\mathrm{Low}\times\mathrm{Post}$ coefficient from a placebo
DiD that assigns a fake reform to the indicated quarter, estimated using only pre-reform
data (2021Q1--2022Q1). Standard errors clustered by department in parentheses; all estimates
are weighted by ENAHO survey weights. $^{*}p<0.1$, $^{**}p<0.05$, $^{***}p<0.01$.
\end{tablenotes}
\end{threeparttable}
\end{table}
